\section{The compactification and its boundary lattice}\label{sec:surface}

The geometry will first be described over an arbitrary field containing
\(s,t\ne0\). Intersection and component arguments will then be made
over its algebraic closure. All blowup centres and formulas are defined
over the ground field.

\subsection{Eight blowups and exactly four accessible lines}

Start with \(\PP^1_x\times\PP^1_y\), and make the eight blowups
in Table~\ref{tab:blowups}. Centres over disjoint points may be treated
in either order. Let \(S=S_{t,s}\) be the resulting smooth projective
rational surface. Write \(H_x,H_y\) for the pullbacks of the classes
of \(x=\text{constant}\) and \(y=\text{constant}\). The symbols
\(E_i\) denote total exceptional classes, rather than strict
transforms throughout the sequence.

\begin{table}[htbp]
\centering
\caption{The eight blowups of \(\PP^1_x\times\PP^1_y\).
Each \(E_i\) denotes the total exceptional class. The conditions for
infinitely near centres are evaluated on the indicated exceptional curve.}
\label{tab:blowups}
\begin{tabularx}{\textwidth}{@{}lX@{}}
\toprule
Class & Centre \\
\midrule
\(E_1\) & \((x,y)=(\infty,1)\). \\
\(E_2\) & \((x,y)=(0,\infty)\). \\
\(E_3\) & On the first exceptional curve over \((0,\infty)\), at \(xy=t\). \\
\(E_4\) & \((x,y)=(0,0)\). \\
\(E_5\) & The intersection of the first exceptional curve over \((0,0)\)
with the strict transform of \(y=0\). \\
\(E_6\) & On the exceptional curve created by the preceding blowup, at \(x^2/y=t\). \\
\(E_7\) & \((x,y)=(\infty,\infty)\). \\
\(E_8\) & On the first exceptional curve over \((\infty,\infty)\), at \(y/x=s\). \\
\bottomrule
\end{tabularx}
\end{table}


For the three infinitely near smooth-boundary centres \(E_3,E_6,E_8\),
coordinates immediately before the last blowup are respectively
\begin{equation}\label{eq:pre-blowup-charts}
(x,y)=(uv,u^{-1}),\qquad
(uv,u^2v),\qquad ((uv)^{-1},u^{-1}).
\end{equation}
The centres are \((u,v)=(0,t),(0,t),(0,s)\). The remaining
smooth-boundary blowup \(E_1\) has centre \(x^{-1}=0,y=1\).
Thus none of these four centres is a boundary node, in any
characteristic, because \(s,t\ne0\).

The free Picard basis is \(H_x,H_y,E_1,\ldots,E_8\), with
intersection form
\begin{equation}\label{eq:intersection-form}
H_x^2=H_y^2=0,\quad H_xH_y=1,\quad
E_iE_j=-\delta_{ij},\quad H_xE_i=H_yE_i=0.
\end{equation}
The following classes are the actual irreducible boundary curves:
\begin{equation}\label{eq:boundary-classes}
\begin{array}{ll}
D_1=H_x-E_1-E_7,&D_2=E_7-E_8,\\
D_3=H_y-E_2-E_7,&D_4=E_2-E_3,\\
D_5=H_x-E_2-E_4,&D_6=E_4-E_5,\\
D_7=E_5-E_6,&D_8=H_y-E_4-E_5.
\end{array}
\end{equation}
They form a transverse cycle in the displayed order. Each has square
\(-2\); consecutive components meet once, including \(D_8,D_1\);
all other intersections vanish. One can see this directly by making
the four corner blowups \(E_2,E_4,E_5,E_7\), followed by the four
smooth-boundary blowups. It also follows from
\eqref{eq:intersection-form} and \eqref{eq:boundary-classes}.
Their sum is
\begin{equation}\label{eq:anticanonical-cycle}
D=\sum_{i=1}^8D_i
=2H_x+2H_y-\sum_{i=1}^8E_i=-K_S,
\qquad D^2=0,\quad DD_i=0.
\end{equation}

Set \(U=U_{t,s}=S\setminus D\). Its points are the disjoint union
of the torus and the four accessible exceptional curves
\(E_1,E_3,E_6,E_8\), each with its one boundary point removed.
Those four curves are affine lines. Neighbourhood coordinates are
\begin{equation}\label{eq:accessible-charts}
\begin{array}{lll}
L_1:&x=u^{-1},&y=1+uv,\\
L_2:&x=u(t+uv),&y=u^{-1},\\
L_3:&x=u(t+uv),&y=u^2(t+uv),\\
L_4:&x=[u(s+uv)]^{-1},&y=u^{-1}.
\end{array}
\end{equation}
In each row the line is \(u=0\), with arbitrary finite coordinate
\(v\). These are the source's accessible charts
\citep[Section 2.1]{JoshiRoffelsen2026}. Their point decomposition is
exactly the state set in Section~\ref{sec:states}. The intermediate
exceptional curves occur in \(D\) and are not additional ordinary
states.

\subsection{Regularity of every forward branch}

The torus formula \eqref{eq:torus-map-intro} extends to a morphism
\(F_t:U_{t,s}\to U_{st,s}\). We give the local formulas, since
regularity at the exceptional points will later extend the invariant.
In each target chart use coordinates \((a,b)\) and phase \(T=st\)
in place of \((u,v)\) and \(t\) in \eqref{eq:accessible-charts}.

Near the torus divisor \(sx-y=0\), target \(L_1\) coordinates are
\begin{equation}\label{eq:extension-torus}
a=\frac{sx-y}{st},\qquad b=\frac{st}{y}.
\end{equation}
They give \(X=a^{-1}\) and \(Y=1+ab\) where \(a\ne0\),
and extend to \(a=0\).

On input \(L_1\), put \(A=1+uv\) and \(B=s-uA\).
Target \(L_2\) coordinates are
\begin{equation}\label{eq:extension-L1}
a=\frac{uA}{s},\qquad
b=\frac{s^2t(A^2-sv)}{A^2B}.
\end{equation}
On input \(L_2\), put \(A=t+uv\) and \(B=su^2A-1\).
Target \(L_3\) coordinates are
\begin{equation}\label{eq:extension-L2}
a=\frac{uAB}{t},\qquad
b=\frac{st^2(-v+2suA^2-s^2u^3A^3)}{A^2B^3}.
\end{equation}
On input \(L_3\), put \(A=t+uv\). Target \(L_4\)
coordinates are
\begin{equation}\label{eq:extension-L3}
a=\frac us,\qquad b=\frac{s(sv-A)}t.
\end{equation}
Finally, on input \(L_4\), put \(A=s+uv\). Target \(L_1\)
coordinates are
\begin{equation}\label{eq:extension-L4}
a=-\frac{v}{stA},\qquad b=stu.
\end{equation}

These are rational identities obtained by inserting
\eqref{eq:accessible-charts} into \eqref{eq:torus-map-intro}.
For example, in \eqref{eq:extension-L2} the identity
\[
t-AB^2=u\bigl(-v+2suA^2-s^2u^3A^3\bigr)
\]
gives the second coordinate after cancelling the displayed factor
\(u\). At a relevant input-line point, every denominator in
\eqref{eq:extension-L1}--\eqref{eq:extension-L4} is a unit:
the values of \(A,B\) are among \(1,s,t,-1\).
The formulas therefore define morphisms on neighbourhoods of every
line point. No division by \(2\) or \(3\) occurs.

Evaluating them at \(u=0\) gives the branches in
Table~\ref{tab:native}. In \eqref{eq:extension-L4}, a nonzero
\(v\) gives the torus point \((-s^2t/v,1)\), whereas \(v=0\)
gives the target \(L_1\) point \((a,b)=(0,0)\). Thus that
special point is included by a regular local map. Together with the
regular torus branch these formulae cover \(U\), and agreement on
the dense torus shows that they define the claimed morphism.

\subsection{Curves orthogonal to the boundary}

The next lemma is the lattice input to the fibre argument. Work over
\(\kkbar\) for the remainder of this section.

\begin{lemma}\label{lem:boundary-lattice}
Every integral curve \(Z\) on \(S\) disjoint from \(D\) is
linearly equivalent to \(dD\) for a positive integer \(d\).
\end{lemma}

\begin{proof}
Write the integral divisor class in the free Picard basis as
\[
Z=aH_x+bH_y-\sum_{i=1}^{8}m_iE_i.
\]
Here the coefficients are integers, not assumed to be multiplicities
of a plane-curve presentation. The eight equations \(ZD_i=0\)
are
\begin{align*}
b-m_1-m_7&=0,& m_7-m_8&=0,\\
a-m_2-m_7&=0,& m_2-m_3&=0,\\
b-m_2-m_4&=0,& m_4-m_5&=0,\\
m_5-m_6&=0,& a-m_4-m_5&=0.
\end{align*}
Their full integer solution has two parameters \(u,v\in\ZZ\):
\begin{equation}\label{eq:lattice-solution}
\begin{split}
a&=2u,\qquad b=u+v,\qquad m_1=2v-u,\\
m_2=m_3&=v,\qquad m_4=m_5=m_6=u,\qquad
m_7=m_8=2u-v.
\end{split}
\end{equation}
Substitution in the intersection form gives
\[
Z^2=2ab-\sum m_i^2=-8(u-v)^2.
\]
Since \(K_S=-D\), adjunction for the integral curve \(Z\) yields
\[
\pa(Z)=1+\frac{Z^2+ZK_S}{2}=1-4(u-v)^2.
\]
An integral proper curve over an algebraically closed field has
\(\pa(Z)=\dim H^1(Z,\OO_Z)\ge0\). Therefore \(u=v=d\),
and \eqref{eq:lattice-solution} becomes the class \(dD\).
The calculation takes place in the actual free Picard group, so it
gives linear equivalence, not only numerical equivalence.
Intersecting with an ample divisor shows \(d>0\), since both
\(Z\) and \(D\) are nonzero effective divisors.
\end{proof}

The intersection form and adjunction compute integers. In particular,
the coefficients \(8\) and \(4\) in this argument are not reduced
modulo the characteristic of the ground field.
